%! AP Statistics — Unit 1, Lesson 9
\documentclass[10pt]{article}
\usepackage[margin=0.6in]{geometry}
\usepackage{xcolor}
\usepackage[most]{tcolorbox}
\usepackage{enumitem}
\usepackage{multicol}
\usepackage{amsmath,amssymb}
\usepackage{array}
\usepackage{tabularx}
\tcbuselibrary{breakable}

\definecolor{sky}{HTML}{EAF3FB}
\definecolor{navy}{HTML}{1F3A5F}
\definecolor{redbox}{HTML}{FCECEC}
\definecolor{greenbox}{HTML}{EDF8F0}
\definecolor{goldbox}{HTML}{FFF7E6}
\newtcolorbox{skillbox}[2][]{
    colback=#2, colframe=navy, boxrule=0.8pt, arc=2mm,
    left=2mm,right=2mm,top=1.5mm,bottom=1.5mm,
    fonttitle=\bfseries, title={#1}, halign title=center, breakable
}
\setlist[itemize]{leftmargin=1.2em,itemsep=0.35em,topsep=0.35em}
\setlist[enumerate]{leftmargin=1.4em,itemsep=0.35em,topsep=0.35em}
\newcolumntype{C}[1]{>{\centering\arraybackslash}p{#1}}
\newcolumntype{Y}{>{\centering\arraybackslash}X}

\begin{document}
\def\CourseName{AP Statistics}
\def\SchoolYear{2026--2027}
\def\MeetingLength{55 minutes}
\def\UnitNumberName{Unit 1: Exploring One-Variable Data \quad}
\def\LessonNumberName{Lesson 1.9: Comparing Distributions of a Quantitative Variable}

\begin{center}
    {\Large\bfseries \CourseName: \SchoolYear\ (\MeetingLength) \\
    {\small\bfseries \UnitNumberName \LessonNumberName}}
\end{center}

\begin{tcolorbox}[colback=sky,colframe=navy,boxrule=0.9pt,arc=2mm,
    left=2mm,right=2mm,top=1.5mm,bottom=1.5mm]
    \textbf{Primary Objective:} Students will compare the distributions of a quantitative variable
    for two or more groups using parallel boxplots, back-to-back stemplots, or histograms. Students
    will write comparative statements using SOCS that include context and explicit comparisons of
    center and spread (Big Idea: UNC; AP Skills 2 \& 3).
\end{tcolorbox}

\begin{skillbox}[Priority Ideas \& Skills]{goldbox}
\begin{minipage}[t]{0.48\linewidth}
    \textbf{AP Skills 2 \& 3 --- Data Analysis \& Using Probability and Simulation}
    \begin{itemize}
        \item Construct parallel boxplots and back-to-back stemplots on the same scale.
        \item Write AP-quality comparative statements: use the word ``compare'' (or equivalent); reference both groups; include context.
        \item Compare centers, spreads, shapes, and outliers across groups explicitly.
    \end{itemize}
\end{minipage}
\hfill
\begin{minipage}[t]{0.48\linewidth}
    \textbf{Key Understandings}
    \begin{itemize}
        \item Comparison requires the same scale for both displays.
        \item AP scoring requires \textit{explicit comparative language}: ``Group A has a higher median than Group B'' earns credit; ``Group A's median is 30'' does not by itself.
        \item A difference in centers with similar spreads is stronger evidence of a real difference than overlapping distributions.
    \end{itemize}
\end{minipage}
\end{skillbox}

\begin{skillbox}[{Vocabulary, Concepts, \& Theorems}]{greenbox}
\begin{minipage}[t]{0.48\linewidth}
    \begin{itemize}
        \item \textbf{Parallel boxplots} (side-by-side): two or more boxplots on the same number line; best for comparing medians and spreads across groups
        \item \textbf{Back-to-back stemplot}: shared stem column; leaves for Group 1 on left, Group 2 on right; best for small datasets ($n < 30$)
        \item \textbf{Comparative statement}: a sentence that explicitly names both groups and uses a comparison word (higher, lower, more variable, similar, etc.)
    \end{itemize}
\end{minipage}
\hfill
\begin{minipage}[t]{0.48\linewidth}
    \begin{itemize}
        \item \textbf{Overlapping distributions}: a large overlap suggests groups are similar; small overlap suggests a meaningful difference
        \item \textbf{Variability (spread) comparison}: if one group has a larger IQR or SD, values in that group are more spread out --- less consistent
        \item \textbf{Context}: every comparison must reference what the data represent (not just numbers)
    \end{itemize}
\end{minipage}
\end{skillbox}

\begin{skillbox}[Activate Prior Knowledge \& Spiral Review]{sky}
\begin{minipage}[t]{0.48\linewidth}
    \textbf{Quick Review (3 min)}\\
    Display two individual boxplots (on separate number lines) of the same variable for two groups. Ask: ``Can you compare them? What makes this hard?'' Answer: they need to be on the \textit{same} scale. Today we fix that.
\end{minipage}
\hfill
\begin{minipage}[t]{0.48\linewidth}
    \textbf{Spiral Connections}
    \begin{itemize}
        \item SOCS framework (Lesson 1.6) and five-number summary (Lessons 1.7--1.8) provide the vocabulary for comparisons.
        \item Comparing two groups is a preview of two-sample inference in Units 6--7.
        \item Unit 2: comparing two quantitative variables uses scatterplots rather than parallel boxplots.
    \end{itemize}
\end{minipage}
\end{skillbox}

\begin{skillbox}[Hook \& Lesson]{sky}
\begin{minipage}[t]{0.48\linewidth}
    \textbf{Hook (5 min)}\\
    Present data: ACT scores for students who did vs.\ did not participate in a test-prep course. Ask: ``Did the test prep course help? How would you decide?'' Students quickly realize they need to compare the two groups on the same display. This motivates parallel boxplots and comparative language.
\end{minipage}
\hfill
\begin{minipage}[t]{0.48\linewidth}
    \textbf{Lesson Flow (15 min)}
    \begin{enumerate}
        \item Construct parallel boxplots of the ACT data on one number line.
        \item Model writing a full AP comparison: shape, outliers, center (with explicit comparison), spread (with explicit comparison).
        \item Discuss: ``Is the difference in medians large relative to the spread? What does overlap tell us?''
        \item Construct a back-to-back stemplot for a smaller dataset; compare with the boxplot version.
    \end{enumerate}
\end{minipage}
\end{skillbox}

\begin{skillbox}[Explicit Instruction \& Active Monitoring]{sky}
\begin{minipage}[t]{0.48\linewidth}
    \textbf{Direct Instruction (10 min)}
    \begin{itemize}
        \item Show an AP model response and label which sentences earn points (comparative language + context + both groups named).
        \item Non-example: ``Group A's median is 45.'' Example: ``Group A's median score (45) was 8 points higher than Group B's median (37).''
        \item Emphasize: always compare center AND spread --- a large difference in center with large spread may not be meaningful.
    \end{itemize}
\end{minipage}
\hfill
\begin{minipage}[t]{0.48\linewidth}
    \textbf{Active Monitoring}
    \begin{itemize}
        \item As students write comparisons, check: Is both groups' context named? Is a comparison word used? Are both center and spread addressed?
        \item Common error: listing statistics for each group separately rather than comparing.
        \item Prompt: ``Don't just describe each group. Tell me which is \textit{bigger}, \textit{more spread}, \textit{more skewed}.''
    \end{itemize}
\end{minipage}
\end{skillbox}

\begin{skillbox}[Group Work \& Differentiation]{redbox}
\begin{minipage}[t]{0.48\linewidth}
    \textbf{Group Activity (8--10 min)}\\
    Groups receive data on hours of sleep for athletes vs.\ non-athletes. Groups:
    \begin{enumerate}
        \item Construct parallel boxplots on the same number line.
        \item Write a complete AP comparison (must include at least 4 explicit comparative statements).
        \item Peer-review another group's response: Does it use comparative language? Context? Cover SOCS?
    \end{enumerate}
\end{minipage}
\hfill
\begin{minipage}[t]{0.48\linewidth}
    \textbf{Differentiation}
    \begin{itemize}
        \item \textit{Support}: Provide a sentence frame: ``The median [variable] for [Group A] was [value], which is [higher/lower] than the median for [Group B] ([value]).'
        \item \textit{Extension}: Compare three groups. Discuss how conclusions become more complex and what additional patterns emerge.
        \item \textit{AP Practice}: Use a released AP free-response prompt requiring comparison; score using the AP rubric.
    \end{itemize}
\end{minipage}
\end{skillbox}

\begin{skillbox}[Individual Work \& Assessment]{redbox}
\begin{minipage}[t]{0.48\linewidth}
    \textbf{Exit Ticket (5 min)}\\
    Given parallel boxplots of monthly rainfall for two cities, students write a 3--4 sentence comparison that includes: center (with explicit comparison), spread (with explicit comparison), and any notable features (outliers, skewness), all with context.
\end{minipage}
\hfill
\begin{minipage}[t]{0.48\linewidth}
    \textbf{AP-Style Check}\\
    \textit{``Compare the distributions of [variable] for the two groups shown in the parallel boxplots. Justify your conclusion.''}\\[4pt]
    Scoring focus: (1) comparative language used, (2) both groups identified with context, (3) center compared, (4) spread compared, (5) other features noted (outliers, shape).
\end{minipage}
\end{skillbox}

\begin{skillbox}[Reinforcement \& Extension]{goldbox}
\begin{minipage}[t]{0.48\linewidth}
    \textbf{Homework / Practice}
    \begin{itemize}
        \item Using a provided dataset (two groups, 15 values each), construct parallel boxplots and write a complete AP comparison.
        \item Practice: AP released FRQ involving comparison of two distributions.
    \end{itemize}
\end{minipage}
\hfill
\begin{minipage}[t]{0.48\linewidth}
    \textbf{Extension / Preview}
    \begin{itemize}
        \item \textit{Extension}: Find a published comparison of two groups (e.g., a news article) and critique the statistical language used. Is it rigorous? What is missing?
        \item \textit{Preview}: Lesson 1.10 introduces the Normal distribution --- a specific symmetric, bell-shaped model that allows us to calculate precise proportions using $z$-scores.
    \end{itemize}
\end{minipage}
\end{skillbox}

\end{document}
